% Options for packages loaded elsewhere
\PassOptionsToPackage{unicode}{hyperref}
\PassOptionsToPackage{hyphens}{url}
\documentclass[
]{article}
\usepackage{xcolor}
\usepackage[margin=1in]{geometry}
\usepackage{amsmath,amssymb}
\setcounter{secnumdepth}{-\maxdimen} % remove section numbering
\usepackage{iftex}
\ifPDFTeX
  \usepackage[T1]{fontenc}
  \usepackage[utf8]{inputenc}
  \usepackage{textcomp} % provide euro and other symbols
\else % if luatex or xetex
  \usepackage{unicode-math} % this also loads fontspec
  \defaultfontfeatures{Scale=MatchLowercase}
  \defaultfontfeatures[\rmfamily]{Ligatures=TeX,Scale=1}
\fi
\usepackage{lmodern}
\ifPDFTeX\else
  % xetex/luatex font selection
\fi
% Use upquote if available, for straight quotes in verbatim environments
\IfFileExists{upquote.sty}{\usepackage{upquote}}{}
\IfFileExists{microtype.sty}{% use microtype if available
  \usepackage[]{microtype}
  \UseMicrotypeSet[protrusion]{basicmath} % disable protrusion for tt fonts
}{}
\makeatletter
\@ifundefined{KOMAClassName}{% if non-KOMA class
  \IfFileExists{parskip.sty}{%
    \usepackage{parskip}
  }{% else
    \setlength{\parindent}{0pt}
    \setlength{\parskip}{6pt plus 2pt minus 1pt}}
}{% if KOMA class
  \KOMAoptions{parskip=half}}
\makeatother
\usepackage{graphicx}
\makeatletter
\newsavebox\pandoc@box
\newcommand*\pandocbounded[1]{% scales image to fit in text height/width
  \sbox\pandoc@box{#1}%
  \Gscale@div\@tempa{\textheight}{\dimexpr\ht\pandoc@box+\dp\pandoc@box\relax}%
  \Gscale@div\@tempb{\linewidth}{\wd\pandoc@box}%
  \ifdim\@tempb\p@<\@tempa\p@\let\@tempa\@tempb\fi% select the smaller of both
  \ifdim\@tempa\p@<\p@\scalebox{\@tempa}{\usebox\pandoc@box}%
  \else\usebox{\pandoc@box}%
  \fi%
}
% Set default figure placement to htbp
\def\fps@figure{htbp}
\makeatother
\setlength{\emergencystretch}{3em} % prevent overfull lines
\providecommand{\tightlist}{%
  \setlength{\itemsep}{0pt}\setlength{\parskip}{0pt}}
\usepackage{graphicx}
\usepackage{float}
\usepackage{bookmark}
\IfFileExists{xurl.sty}{\usepackage{xurl}}{} % add URL line breaks if available
\urlstyle{same}
\hypersetup{
  hidelinks,
  pdfcreator={LaTeX via pandoc}}

\author{}
\date{\vspace{-2.5em}}

\begin{document}

\section{Question}\label{question}

Un experimento consiste en medir la elongación de un resorte en función
de la fuerza aplicada. En la gráfica se registran los resultados de 93
mediciones realizadas con el mismo resorte.

\includegraphics[width=12cm,height=\textheight,keepaspectratio]{dispersion_resorte.png}

El comportamiento de la elongación respecto a la fuerza aplicada es

\subsection{Answerlist}\label{answerlist}

\begin{itemize}
\tightlist
\item
  proporcional y mas variable cuanto mayor sea la carga.
\item
  proporcional y mas variable cuanto mayor sea la elongacion.
\item
  de tipo no lineal y mas variable cuanto mayor sea la carga.
\item
  de tipo no lineal y mas variable cuanto mayor sea la elongacion.
\end{itemize}

\section{Solution}\label{solution}

Para analizar el comportamiento de la elongación respecto a la fuerza
aplicada, debemos observar dos aspectos clave en la gráfica de
dispersión:

\textbf{1. Tipo de relación (proporcional vs de tipo no lineal):}

La relación entre la fuerza y la elongación sigue la Ley de Hooke:

\[x = \frac{F}{k}\]

donde \(x\) es la elongación, \(F\) es la fuerza aplicada y \(k\) es la
constante del resorte. Esta es una función \textbf{proporcional}
(proporcionalidad directa). En la gráfica se observa claramente una
tendencia rectilínea ascendente.

\textbf{2. Patrón de dispersión:}

Observando la gráfica:

\begin{itemize}
\tightlist
\item
  En fuerzas pequeñas (cercanas a 0-2 N): los puntos están más
  concentrados
\item
  En fuerzas grandes (cercanas a 10-12 N): los puntos muestran mayor
  variabilidad
\end{itemize}

Esto indica que \textbf{la dispersión es proporcional a la carga}, no a
la elongacion. Físicamente, esto se explica porque la incertidumbre en
la aplicación de fuerzas mayores es más difícil de controlar, y pequeñas
variaciones en la fuerza aplicada producen mayores fluctuaciones en la
medición.

\textbf{Conclusión:}

El comportamiento es \textbf{proporcional} (Ley de Hooke) y \textbf{mas
variable cuanto mayor sea la carga} (mayor variabilidad en fuerzas
grandes).

\subsection{Answerlist}\label{answerlist-1}

\begin{itemize}
\tightlist
\item
  \textbf{Verdadero}. La relación es proporcional (Ley de Hooke) y la
  dispersión aumenta proporcionalmente a la carga.
\item
  \textbf{Falso}. Aunque la relación es proporcional, la dispersión NO
  aumenta con la elongacion sino con la carga.
\item
  \textbf{Falso}. La relación NO es de tipo no lineal (es proporcional
  según la Ley de Hooke).
\item
  \textbf{Falso}. La relación NO es de tipo no lineal (es proporcional)
  y la dispersión NO depende de la elongacion.
\end{itemize}

\section{Meta-information}\label{meta-information}

exname:
elongacion\_resorte\_ley\_hooke\_interpretacion\_representacion\_n2\_v1
extype: schoice exsolution: 1000 exshuffle: TRUE exsection:
Estadistica/Graficas de Dispersion

\end{document}
